% Generated draft: requires mathematical and source review.
\documentclass[11pt]{article}
\usepackage[margin=25mm]{geometry}
\usepackage{fontspec}
\setmainfont{Latin Modern Roman}
\usepackage{amsmath,amssymb,mathtools,unicode-math}
\setmathfont{Latin Modern Math}
\usepackage{enumitem,needspace,xurl,hyperref}
\setlength{\parindent}{0pt}
\setlength{\parskip}{5pt}
\setlength{\emergencystretch}{3em}
\newcommand{\N}{\mathbb{N}}
\newcommand{\Z}{\mathbb{Z}}
\newcommand{\Q}{\mathbb{Q}}
\newcommand{\R}{\mathbb{R}}
\newcommand{\C}{\mathbb{C}}
\newcommand{\F}{\mathbb{F}}
\newcommand{\A}{\mathbb{A}}
\newcommand{\PP}{\mathbb{P}}
\newcommand{\E}{\mathbb{E}}
\newcommand{\eps}{\varepsilon}
\newcommand{\dd}{\,\mathrm{d}}
\begin{document}
\section*{UnsolvedMath ID: 800136}
\textbf{Source problem code:} catalogue entry
\subsection{Definitions and mathematical statement}
The supplied record combines three contributions under Problem 5 of the cited problem list. It does not specify a single common conjecture.

\textbf{Nathan Ng's discrete moment questions.}
Let \(\zeta(s)\) denote the Riemann zeta function, and write a nontrivial zero as \(\rho=\beta+i\gamma\), where \(0<\beta<1\) and \(\gamma\in\R\). For \(T>0\) and \(k>0\), the source considers the expression
\[
M_k(T)=\sum_{\substack{\rho:\,\zeta(\rho)=0\\0<\operatorname{Im}\rho<T}}
\left|\frac{\zeta(2\rho)}{\zeta'(\rho)}\right|^{2k}.
\]
This expression requires a qualification absent from the supplied statement: division by \(\zeta'(\rho)\) requires the zeros in the sum to be simple. No convention for multiple zeros, no zero-counting convention, and no Riemann hypothesis assumption are supplied. Thus the displayed expression is a proposed moment, defined as written when all included zeros are simple; this draft does not add simplicity as an established fact.

The source attributes to Ng's Ph.D. thesis [NgPhd] the conjecture
\[
M_1(T)\sim\frac{T}{2\pi}\qquad(T\to\infty).
\]
Its first requested task is to use the argument of Milinovich and Ng [MN] to establish
\[
M_1(T)\geq\left(\frac{1}{4\pi}+o(1)\right)T
\qquad(T\to\infty).
\]
Whenever the moment is defined, this lower-bound target can equivalently be expressed as
\[
\liminf_{T\to\infty}\frac{M_1(T)}{T}\geq\frac{1}{4\pi}.
\]
Its second requested task is to prove lower bounds of the displayed shape
\[
M_k(T)\gg T(\log T)^{(k+1)^2}\qquad(T\to\infty).
\]
The range of \(k\) for this task is not explicitly specified. The immediately following comparison concerns \(k>0\), but that does not resolve the intended range or exponent in this task. For a fixed admissible \(k\), the notation \(\gg\) means a lower bound by a positive constant times the right-hand side for all sufficiently large \(T\), with the constant permitted to depend on \(k\).

There is an unresolved incompatibility: if the second task includes \(k=1\), it requests \(M_1(T)\gg T(\log T)^4\), which cannot hold together with the preceding conjecture \(M_1(T)\sim T/(2\pi)\). Both formulas are retained exactly as supplied; neither the exponent nor the range of \(k\) is corrected here.

For comparison, the source attributes to Heap, Li, and Zhao [HLZ] and Gao and Zhao [GZ] lower bounds of the form
\[
\sum_{0<\gamma<T}\frac{1}{|\zeta'(\rho)|^{2k}}
\gg T(\log T)^{(k+1)^2},\qquad k>0.
\]
This is a supplied literature claim, not an independently verified theorem in this draft. Its hypotheses, treatment of multiple zeros, and precise summation convention are not provided.

\textbf{Nicol Leong's multiplicity question.}
For a nontrivial zero \(\rho\), let \(m(\rho)\) be its order of vanishing: the integer \(m\geq1\) such that \(\zeta^{(j)}(\rho)=0\) for \(0\leq j<m\) and \(\zeta^{(m)}(\rho)\neq0\). The question is to improve upper bounds for \(m(\beta+i\gamma)\) as \(\gamma\to+\infty\). The source reports the unconditional bound
\[
m(\beta+i\gamma)\ll\log\gamma,
\]
and, respectively under the Lindel\"of hypothesis and the Riemann hypothesis, the bounds
\[
m(\beta+i\gamma)=o(\log\gamma),
\qquad
m(\beta+i\gamma)\ll\frac{\log\gamma}{\log\log\gamma}.
\]
Here the Riemann hypothesis states that every nontrivial zero has real part \(1/2\), and the Lindel\"of hypothesis states that, for every \(\varepsilon>0\),
\[
|\zeta(1/2+it)|\ll_{\varepsilon}(1+|t|)^{\varepsilon}
\qquad(|t|\to\infty).
\]
The reported bounds concern all nontrivial zeros of sufficiently large positive ordinate. A bound \(\ll\) has an implied constant independent of the zero; the little-oh bound is understood uniformly over such zeros.

The source attributes further unconditional results of order \(\ll\log\gamma\) to Ivi\'c [Ivic1, Ivic2], and attributes to Karatsuba the conjectured bound
\[
m(\beta+i\gamma)\ll1.
\]
The latter asks for a constant bound independent of the zero. The source also states that its ``LI conjecture'' would imply \(m(\rho)=1\) for every nontrivial zero. Since LI is not defined in the record, that implication is preserved as a source qualification without assigning it a precise hypothesis here. The requested improvement has no specified quantitative target; Karatsuba's bound is presented as a conjectured benchmark.

\textbf{Alexandre Bailleul's question about bias and first reversals.}
For real \(x>1\), define
\[
\pi(x)=\#\{p\leq x:p\text{ is prime}\},
\qquad
\operatorname{li}(x)=\operatorname{PV}\int_0^x\frac{\mathrm{d}t}{\log t}.
\]
For an integer \(q\geq2\) and an integer \(a\), define
\[
\pi(x;q,a)=\#\{p\leq x:p\text{ is prime and }p\equiv a\pmod q\}.
\]
For a measurable set \(E\subseteq[2,\infty)\), its logarithmic density, when the limit exists, is
\[
\delta(E)=\lim_{X\to\infty}\frac{1}{\log X}
\int_2^X\mathbf{1}_E(t)\,\frac{\mathrm{d}t}{t}.
\]
These definitions describe the quantities used in the source; they do not resolve its missing hypotheses for density existence.

The source calls the smallest \(x\) with \(\pi(x)>\operatorname{li}(x)\) ``Skewes's number.'' The search domain must be specified for that description to be precise. With the principal-value definition above, an unrestricted search over \(x>1\) includes small values near \(1\), so the intended large first crossing requires an additional domain restriction. No such restriction is stated in the record.

For the intended large crossing, the source credits Littlewood with existence and Skewes with an initial enormous bound, and reports the interval
\[
\left(10^{19},\;1.3971653\times10^{316}\right).
\]
It attributes the lower endpoint bound to B\"uthe [Buthe0] and the upper endpoint bound to Saouter, Trudgian, and Demichel [STD]. It says that the value remains unknown and that experts suggest a value close to the upper bound; the latter is an unattributed expectation, not a quantitative claim.

The source attributes to Rubinstein and Sarnak [RS] the approximate logarithmic density
\[
\delta\bigl(\{x\geq2:\pi(x)>\operatorname{li}(x)\}\bigr)
\approx2.6\times10^{-7}.
\]
For the race between primes congruent to \(2\pmod3\) and primes congruent to \(1\pmod3\), it reports a corresponding density, denoted there by \(\delta_{q;1,2}\), of approximately \(1.0\times10^{-3}\), again attributed to [RS]. The record does not explicitly identify the inequality represented by that notation. It attributes to Bays and Hudson the assertion that the first sign change of
\[
\pi(x;3,2)-\pi(x;3,1)
\]
occurs at \(x=608{,}981{,}813{,}029\). The treatment of ties and the meaning of ``first sign change'' are not specified.

The proposed research direction is to investigate whether the size of the first reversal in a prime number race is correlated with the strength of its logarithmic-density bias. The source further says that classical prime number races have logarithmic densities tending to \(1/2\) as \(q\to\infty\), and suggests that their corresponding first-reversal values should become less extreme, while remaining bounded below by \(q\). The family of races, hypotheses, normalization of first-reversal size, meaning of correlation, and conditions for the asserted lower bound by \(q\) are unspecified. Consequently this contribution supplies a heuristic research question, not a fully quantified conjecture.
\subsection{Short English statement}
Permanent UnsolvedMath ID \(800136\) bundles questions about discrete zeta moments, multiplicities of zeta zeros, and the relationship between prime-race bias and first reversals. The supplied moment exponent conflicts with the supplied first-moment conjecture if \(k=1\) is included. Simplicity assumptions, several literature hypotheses, and a precise formulation of the prime-race question remain unresolved. This is a statement-only draft requiring human mathematical review.
\subsection{Sources}
\begin{itemize}
\item Local UnsolvedMath registry: permanent ID \(800136\), source problem code \texttt{catalogue entry}, title ``Comparative Prime Number Theory --- Problem 5,'' in the AMR Open Problem Lists collection. The registry labels the record \texttt{open}; this is supplied metadata, not an independently verified status assessment. The supplied upstream revision is \texttt{d371ee917ac9647fbb06cd6fc529deec6f3dda82}.
\item Hamieh, Kadiri, Martin, and Ng, \emph{Comparative Prime Number Theory Problem List} (2024), source item Problem (\#5), \url{https://arxiv.org/abs/2407.03530}. The supplied provenance records access on 2026-10-16 and extraction by \texttt{arxiv\_latex}; these are supplied dated provenance, not a claim of independent source access or verification. Registry creation and update are dated 2026-10-16.
\item The individual contributions are attributed to Nathan Ng, University of Lethbridge; Nicol Leong, University of New South Wales--Canberra at ADFA; and Alexandre Bailleul, Universit\'e Paris--Saclay. The registry attributes the collected problem to Hamieh, Kadiri, Martin, and Ng, with proposed year 2024.
\item [NgPhd]: Ng's Ph.D. thesis, cited in the supplied source for the first-moment conjecture. No full bibliographic entry or separate URL is supplied.
\item [MN]: Milinovich and Ng, cited for an argument proposed for the first-moment lower bound. No title, date, full bibliographic entry, or separate URL is supplied.
\item [HLZ]: Heap, Li, and Zhao; [GZ]: Gao and Zhao. The supplied source attributes lower bounds for negative discrete moments of \(\zeta'\) to these articles, without supplying their full bibliographic entries, URLs, or theorem hypotheses.
\item [Ivic1, Ivic2]: works of Ivi\'c, cited for multiplicity bounds and the associated approach. The record supplies no full bibliographic entries or separate URLs. Karatsuba is credited with the conjectured bounded-multiplicity estimate, without a separate citation.
\item [Buthe0]: B\"uthe, credited with the stated lower bound for the intended first crossing of \(\pi\) and \(\operatorname{li}\). [STD]: Saouter, Trudgian, and Demichel, credited with the stated upper bound. No full bibliographic entries or separate URLs are supplied.
\item [RS]: Rubinstein and Sarnak, credited with the stated approximate logarithmic densities. Bays and Hudson are credited with the stated first sign change for the modulus-\(3\) race. Littlewood and Skewes are credited with the existence result and an initial bound, respectively. No further bibliographic details, separate URLs, or relevant theorem hypotheses are supplied for these claims.
\end{itemize}
\end{document}
